\section{Experiments}
\label{sec:experiments}

Every experiment in this section has a \emph{pre-registered decision
rule} that lives in its own report under \texttt{docs/experiments/}.
The paper never re-derives numbers; it cites the experiment id.

\subsection{RQ1: answer-match vs.\ trace/effect alignment (e4.4)}
\label{sec:rq1}

\paragraph{Setup.} Same 56 forward-distilled specs $\times$ 3 candidate
models (\texttt{anthropic/claude-haiku-3.5},
\texttt{anthropic/claude-haiku-4.5}, \texttt{qwen/qwen3-8b}); same
replay-mode evaluation. Both scorers consume the same eval JSONLs. The
answer-match scorer
(\texttt{dmcp/baselines/answer\_match.py}) is token-Jaccard with
substring fallback; threshold $0.5$.

\paragraph{Decision rule (pre-registered).} Positive iff Kendall's
$\tau$ between the two model rankings $< 0.5$ and false-pass $>$
false-fail.

\paragraph{Result.} $\tau = \mathbf{-0.816}$ (rankings nearly reversed).
The strongest model under trace-align (\texttt{haiku-4.5}, 45\,\%)
becomes the worst under answer-match (89\,\%, tied with \texttt{qwen3});
the weakest under trace-align (\texttt{haiku-3.5}, 21\,\%) becomes the
best under answer-match (95\,\%). Overall false-pass rate is
$\mathbf{63.1\%}$; false-fail rate is $\mathbf{1.8\%}$. The headline
visual is Figure~\ref{fig:rq1_kendall} (shown in \S\ref{sec:intro}).

\paragraph{Implication.} A leaderboard scored by final-answer string
match \emph{would attribute the best score to the model that least
often does the right thing}. The CLAUDE.md ``never grade the final
answer'' invariant acquires a quantitative defense.

\subsection{RQ2: forward vs.\ graph-sampling vs.\ direct generation (e4.3)}
\label{sec:rq2}

\paragraph{Setup.} Three TaskSpec generation paths on a shared
substrate (\texttt{time + wikipedia + arxiv}): the headline forward
distillation (9 specs filtered from v3); a graph-sampling baseline
(\texttt{dmcp/baselines/graph\_sampling.py}, 6 specs across
\texttt{chain} and \texttt{hub} motifs); a direct generate-then-verify
baseline (\texttt{dmcp/baselines/direct\_generation.py}, 9 specs). Both
baselines are clearly labeled (\texttt{distiller\_version="baseline-..."})
and never imported by the headline scoring path.

\paragraph{Decision rule (pre-registered).} Positive iff forward
mean $|eq\_set| > 1.05$ and both baselines mean $|eq\_set| = 1.00$.

\paragraph{Result.} Forward mean $|eq\_set| = \mathbf{1.42}$ (max 2,
singleton rate 58\,\%). Both baselines $= 1.00$ exactly, by
construction. Filter pass rate 100\,\% / 100\,\%; coverage 65\,\% /
44\,\% / 18\,\%; zero marker violations. See
Table~\ref{tab:rq2_comparison}.

% AUTO-GENERATED by paper/regenerate.py -- do not hand-edit.
% id: tab:rq2_comparison
% status: ready
% caption: Forward distillation vs graph-sampling vs direct generate-then-verify, on the offline-derivable axes (mean / max &#124;eq_set&#124;, singleton rate, missing-arg-predicate rate, coverage, filter pass rate, executable-by-construction, ordering density).
\begin{table*}[t]
  \centering
  \small
  \begin{tabular}{lrrrrrr}
    \toprule
    method & n & mean $|eq\_set|$ & singleton & missing arg-pred. & coverage & filter pass \\
    \midrule
    forward & 9 & 1.421 & 57.9\% & 21.1\% & 64.7\% & -- \\
    graph & 6 & 1.000 & 100.0\% & 100.0\% & 44.1\% & 100.0\% \\
    direct & 9 & 1.000 & 100.0\% & 0.0\% & 17.6\% & 100.0\% \\
    \bottomrule
  \end{tabular}
  \caption{Forward distillation vs graph-sampling vs direct generate-then-verify, on the offline-derivable axes (mean / max |eq\_set|, singleton rate, missing-arg-predicate rate, coverage, filter pass rate, executable-by-construction, ordering density).}
  \label{tab:rq2_comparison}
\end{table*}


\paragraph{Implication.} Path-agnostic equivalence sets are the
structural difference between forward distillation and the two closest
prior-art shapes. The gap is what makes effect-based scoring faithful
across multiple valid agent trajectories.

\subsection{RQ3: trace-property failure model (e4.5)}
\label{sec:rq3}

\paragraph{Setup.} Ridge-regularised logistic regression of
pass / fail on the \texttt{ComplexityProfile} + dynamism features
(\texttt{trace\_depth}, \texttt{runtime\_branching},
\texttt{state\_coupling}, \texttt{cross\_server},
\texttt{dynamism\_live}, \texttt{dynamism\_stateful}); per-candidate +
pooled. Pure-Python IRLS, no new dependency. Drop-column permutation
importance.

\paragraph{Decision rule (pre-registered).} Positive iff at least one
feature has $|\beta| > 0.5$ in the pooled fit AND that feature's
drop-loglik loss $> 0.5$ AND it is the top driver in $\geq 2$ of $3$
per-model fits.

\paragraph{Result.} Pooled \texttt{dynamism\_live}
$\beta = \mathbf{-3.110}$, odds ratio $\mathbf{0.045}$, drop-loglik loss
$\mathbf{+6.606}$~--- $\approx 70\%$ of the pooled importance budget.
Top driver in \texttt{haiku-4.5} and \texttt{qwen3-8b}; second behind
\texttt{trace\_depth} in \texttt{haiku-3.5}. Secondary effects:
\texttt{trace\_depth} ($\beta = -0.097$, loss $+1.682$) and
\texttt{cross\_server} ($\beta = -0.757$, loss $+0.871$). See
Table~\ref{tab:rq3_failure_drivers}.

% AUTO-GENERATED by paper/regenerate.py -- do not hand-edit.
% id: tab:rq3_failure_drivers
% status: ready
% caption: Pooled coefficients + odds ratios + drop-loglik importance for the RQ3 failure model.
\begin{table}[t]
  \centering
  \small
  \begin{tabular}{lrrr}
    \toprule
    feature & $\beta$ & odds ratio & drop-loglik loss \\
    \midrule
    \texttt{dynamism\_live} & $-3.110$ & 0.045 & $+6.606$ \\
    \texttt{trace\_depth} & $-0.097$ & 0.908 & $+1.682$ \\
    \texttt{cross\_server} & $-0.757$ & 0.469 & $+0.871$ \\
    \texttt{runtime\_branching} & $-0.328$ & 0.720 & $+0.261$ \\
    \texttt{state\_coupling} & $-0.897$ & 0.408 & $+0.000$ \\
    \texttt{dynamism\_stateful} & $-0.897$ & 0.408 & $+0.000$ \\
    \bottomrule
  \end{tabular}
  \caption{Pooled coefficients + odds ratios + drop-loglik importance for the RQ3 failure model.}
  \label{tab:rq3_failure_drivers}
\end{table}


\paragraph{Implication.} Task \emph{dynamism} is not just a taxonomy
label~--- it is the dominant statistical driver of failure across
models. This is the empirical defense of the project's choice to
stratify by dynamism class.

\subsection{RQ4: scorer-vs-human validation (e4.6, status: planned)}
\label{sec:rq4}

\paragraph{Setup.} A deterministic stratified subset of 200 tasks
(\texttt{dmcp rq4-subset}, balanced by dynamism $\times$
complexity\_bin), three human raters per cell, full annotation protocol
pre-registered in
\texttt{docs/experiments/e4.6-rq4-scorer-vs-human.md}. Cohen's $\kappa$
between Tier-1 / Tier-2 verdicts and the majority-vote consensus;
Krippendorff's $\alpha$ over the full rater $\times$ scorer grid;
replay determinism (Tier-1 verdict flip rate between two re-runs).

\paragraph{Decision rule (pre-registered).} Positive iff
$\kappa_{\text{tier1}} \geq 0.70$ AND $\kappa_{\text{tier2}} \geq 0.70$
AND replay flip rate $< 5\%$.

\paragraph{Result.} \emph{Pending the human annotation pass~--- the
harness merged without faking numbers.} The experiment doc captures
the protocol so the result can be filled in deterministically once
raters complete the pass. The placeholder
Table~\ref{tab:rq4_agreement} keeps the cross-reference resolved.

% AUTO-GENERATED by paper/regenerate.py -- do not hand-edit.
% id: tab:rq4_agreement
% status: pending
% reason: human annotation pass not completed yet -- pre-registered in docs/experiments/e4.6-rq4-scorer-vs-human.md; harness ready to consume docs/experiments/e4.6_numbers.json once it lands.
\begin{table}[t]
  \centering
  \fbox{\parbox{0.9\columnwidth}{\centering \textbf{tab:rq4\_agreement} -- placeholder.\\
  Status: pending.\\
  Gating step: E4.6 annotation pass}}
  \caption{Per-tier (Tier-1 / Tier-2) Cohen's $\kappa$ vs human consensus; Krippendorff's $\alpha$ over the full grid; per-tier false-pass / false-fail rates; replay flip rate.}
  \label{tab:rq4_agreement}
\end{table}


\paragraph{Implication (preview).} Whatever the result, it is the
\emph{honest calibration} for every claim in
\S\ref{sec:rq1}--\S\ref{sec:rq3}: if $\kappa \geq 0.70$ the trace-align
scorer is materially aligned with humans; if it is not, the downstream
claims are revisited.

\subsection{Degradation curves: $P_{\text{alt}}$ and complexity stratification}
\label{sec:rq_palt}

The pool sampler (\S\ref{sec:method:controls}) drives a $P_{\text{alt}}$
grid $\in \{0, .25, .5, .75, 1.0\}$ that quantifies how candidate
accuracy and SAE rate degrade as alternative tools dilute the
candidate's tool surface. Results are stratified by complexity bin
($1$ / $2$ / $3$--$4$ / $5+$ required tools) with macro and micro
averages.

% AUTO-GENERATED by paper/regenerate.py -- do not hand-edit.
% id: fig:p_alt_degradation
% status: pending
% reason: no docs/experiments/e2.7_numbers.json committed yet -- the P_alt driver (dmcp curve) is in place but the experiment doc has not been run end-to-end.
\begin{figure}[t]
  \centering
  \fbox{\parbox{0.9\columnwidth}{\centering \textbf{fig:p\_alt\_degradation} -- placeholder.\\
  Status: pending.\\
  Gating step: E2.7 (P\_alt driver, already merged) + an experiment-doc run}}
  \caption{P\_alt degradation curves (accuracy vs P\_alt and SAE rate vs P\_alt) per strategy $\times$ level, with Wilson CIs and complexity-bin facets.}
  \label{fig:p_alt_degradation}
\end{figure}


\subsection{Living bench: decay over time}
\label{sec:rq_decay}

\texttt{dmcp refresh} re-executes each spec's reference trace against
the live substrate at a later time and classifies every call as
identical / drifted / broken / skipped. Reports per-server drift rates
and a decay table; spec staleness is gated by a configurable threshold.

% AUTO-GENERATED by paper/regenerate.py -- do not hand-edit.
% id: fig:decay_curve
% status: pending
% reason: no docs/experiments/e1.5_numbers.json committed yet -- the refresh decay driver is in place but the experiment doc has not been run end-to-end.
\begin{figure}[t]
  \centering
  \fbox{\parbox{0.9\columnwidth}{\centering \textbf{fig:decay\_curve} -- placeholder.\\
  Status: pending.\\
  Gating step: E1.5 (decay metrics + backoff, already merged) + an experiment-doc run on multi-window refresh data}}
  \caption{Per-server decay over time: drift, broken, identical rates as the substrate ages.}
  \label{fig:decay_curve}
\end{figure}


\subsection{Leaderboard and capability profile}
\label{sec:rq_leaderboard}

$\geq 5$ candidate models (a GPT-class, Gemini, Claude Sonnet / Opus,
an open-weight 70B+, a tool-specialised model) under replay,
3$\times$ per task. Per-stratum capability profile $=$ accuracy by
(dynamism, complexity bin, recovery\_required, runtime\_branching).
Pending the substrate scale-up.

% AUTO-GENERATED by paper/regenerate.py -- do not hand-edit.
% id: tab:capability_profile
% status: pending
% reason: >=5-model leaderboard not produced yet -- depends on E4.7 + the larger substrate from E3.1.
\begin{table}[t]
  \centering
  \fbox{\parbox{0.9\columnwidth}{\centering \textbf{tab:capability\_profile} -- placeholder.\\
  Status: pending.\\
  Gating step: E4.7 (leaderboard) --- itself gated on E3.1}}
  \caption{Per-model accuracy stratified by (dynamism $\times$ complexity bin $\times$ recovery\_required $\times$ runtime\_branching); $\geq$ 5 models.}
  \label{tab:capability_profile}
\end{table}


\subsection{Generation-strategy diversity and difficulty}
\label{sec:rq_gen_strategy}

The benchmark spans 15 generation strategies (six base relationship
samplers, seven corner cases, and two cross-server specials) and
difficulty bins ranging from single-tool tasks to long chains of
pairwise-similar tools. Table~\ref{tab:gen_strategy_ablation} reports,
per generation strategy, the agent pass-rate, SAE rate, and
pass\textsuperscript{$k$} reliability --- isolating which strategy
yields the hardest and most attribution-confusing tasks.
Figure~\ref{fig:difficulty_curve} traces the expected monotone
degradation as trace depth grows, and
Figure~\ref{fig:gen_eval_sae_heatmap} crosses generation strategy with
eval condition to locate the combinations that most provoke
server-attribution error.

% AUTO-GENERATED by paper/regenerate.py -- do not hand-edit.
% id: tab:gen_strategy_ablation
% status: pending
% reason: data source `docs/experiments/e4.9_numbers.json` not on disk yet -- the gating step listed in figures.md has not produced it.
\begin{table}[t]
  \centering
  \fbox{\parbox{0.9\columnwidth}{\centering \textbf{tab:gen\_strategy\_ablation} -- placeholder.\\
  Status: pending.\\
  Gating step: E4.9 (gen-strategy ablation) on the diverse corpus}}
  \caption{Per generation strategy (15 strategies: 6 base + 7 corner + 2 special): n, pass-rate, SAE-rate, pass\textasciicircum{}k --- which generation strategy yields the hardest / most-SAE tasks.}
  \label{tab:gen_strategy_ablation}
\end{table}


% AUTO-GENERATED by paper/regenerate.py -- do not hand-edit.
% id: fig:difficulty_curve
% status: pending
% reason: data source `docs/experiments/e4.10_numbers.json` not on disk yet -- the gating step listed in figures.md has not produced it.
\begin{figure}[t]
  \centering
  \fbox{\parbox{0.9\columnwidth}{\centering \textbf{fig:difficulty\_curve} -- placeholder.\\
  Status: pending.\\
  Gating step: E4.10 (difficulty curve) on the diverse corpus}}
  \caption{Agent pass-rate by trace-depth difficulty bin (simple 1--2 / medium 3--4 / hard 5+ long-similar-chain), per candidate model --- the expected monotone degradation.}
  \label{fig:difficulty_curve}
\end{figure}


% AUTO-GENERATED by paper/regenerate.py -- do not hand-edit.
% id: fig:gen_eval_sae_heatmap
% status: pending
% reason: data source `docs/experiments/e4.9_numbers.json` not on disk yet -- the gating step listed in figures.md has not produced it.
\begin{figure}[t]
  \centering
  \fbox{\parbox{0.9\columnwidth}{\centering \textbf{fig:gen\_eval\_sae\_heatmap} -- placeholder.\\
  Status: pending.\\
  Gating step: E4.9 (gen-strategy ablation) on the diverse corpus}}
  \caption{SAE rate per generation strategy (rows) $\times$ eval condition (columns) --- which generation $\times$ which eval setup most provokes server-attribution error.}
  \label{fig:gen_eval_sae_heatmap}
\end{figure}

